\chapter*{Antes de empezar}

Este libro trata de una máquina que pesa kilo y medio, consume veinte vatios, como una bombilla tenue, y aun así reconoce caras, aprende idiomas y atrapa una pelota en el aire. Trata de su cerebro. Lo miraremos como un ingeniero mira una placa desconocida: qué piezas tiene, cómo están conectadas, qué señales corren por los cables y qué calcula ese circuito.

No necesitará fórmulas ni saber biología. Le bastará con curiosidad y con estar dispuesto a imaginar un cable, una bombilla, un interruptor. Todo lo demás lo iremos armando por el camino.

El libro tiene una segunda versión, la completa, con ecuaciones, modelos, ejercicios y referencias a los experimentos. Los capítulos de ambas versiones van en el mismo orden y llevan los mismos números. Si algún capítulo le engancha, abra su versión completa: allí encontrará de dónde sale cada afirmación y cómo comprobarla usted mismo. En el sitio web del libro, cada capítulo se cambia entre “Breve” y “Completo” con un solo botón.

Una advertencia. Del cerebro no se sabe todo, ni mucho menos, y le diré con honestidad dónde termina lo medido y empieza la conjetura. No es un defecto del libro, sino lo más interesante que tiene: muchos capítulos terminan con preguntas que todavía nadie sabe responder.
